\section{Resultados}
\label{subsec:resultados}
A questão norteadora deste trabalho foi: \textbf{Como a adoção de práticas de experimentação contínua pode auxiliar a tomada de decisão sobre a implantação de novas versões de um produto de \textit{software}?}.

Para respondê-la, foi realizado um estudo de caso planejado com base nos insumos levantados por meio de uma revisão da literatura. Dessa questão principal, derivaram-se duas questões específicas para investigar a aplicação das metodologias analisadas e suas implicações: \textbf{Durante o experimento, foi possível identificar diferença na qualidade em uso entre as versões de controle (A) e tratamento (B) do produto a partir de seu uso?} e \textbf{O processo de experimentação contribuiu para uma tomada de decisão baseada em dados?}

Utilizando a metodologia GQM, foi estruturado um protocolo para definir as métricas a serem observadas e analisadas (vide Subseção \ref{subsec:estudo-questao}). A primeira métrica foi o indicador de qualidade em uso, um valor quantitativo agregado de diferentes medidas coletadas do uso do \textit{software} para avaliação da qualidade em uso das versões. A segunda foi a percepção qualitativa dos colaboradores envolvidos no processo, coletada por meio de um questionário semiestruturado.

Inicialmente, foi proposto um modelo de desenvolvimento adaptado ao produto atual para viabilizar um experimento que validasse uma hipótese dos \textit{stakeholders} sobre o \textit{software}, oferecendo uma avaliação quantitativa da nova versão. Esse processo (apresentado no Apêndice \ref{apdx:processo-novo}) baseou-se em trabalhos identificados na revisão da literatura realizada anterior ao planejamento do estudo de caso (Seção \ref{sec:rsl-resultados}).

As primeiras atividades propostas foram relacionadas ao levantamento de suposições relacionadas sobre a plataforma em questão, advindas de conhecimentos de negócio, \textit{feedbacks} de usuários, etc. Das suposições levantadas, a escolhida referia-se ao posicionamento de um elemento na plataforma, pois a equipe percebia que ele estava sendo subutilizado pelos usuários. O propósito então era avaliar a qualidade de uma versão após o reposicionamento comparada à já existente, além de se estudar a adoção da funcionalidade do elemento nos dois casos.

A pesquisa utilizou como referência a norma \citeonline{nbr9241}, que define características e subcaracterísticas da qualidade em uso, e o modelo QATCH \cite{qatch}, que propõe a utilização de técnicas de decisão multicritério para realização do cálculo de um indicador agregado de qualidade. Os dados foram coletados a partir do uso de dois grupos de usuários, cada um exposto consistentemente a uma versão específica do produto. Após a coleta das métricas, foi calculado um indicador para cada usuário único, seguido de um teste de hipótese para verificar a existência de diferença estatística na qualidade em uso percebida entre as populações. Além disso, aplicaram-se técnicas de estatística descritiva para compreender o comportamento dos usuários.

A análise não revelou diferença nos indicadores de qualidade percebida entre as duas versões, no entando, além desta avaliação, a decisão da equipe sobre a implantação da nova versão baseou-se também em métricas individuais, como a adoção da funcionalidade após a mudança de posicionamento e o impacto nas demais funcionalidades afetadas por esse rearranjo na versão de tratamento. De acordo com as métricas de frequência de uso e proporção de adoção, a versão de tratamento mostrou-se benéfica, com a funcionalidade possuindo uma maior proporção de adoção pelos usuários que utilizaram a versão de tratamento sem que as outras fossem afetadas negativamente. Este resultado levou a equipe a optar por sua implementação definitiva.

Um ponto levantado pelos envolvidos após a análise dos resultados foi que essa mudança isolada pode não fornecer insumos suficientes para outras decisões negócio referentes a funcionalidade estudada, tornando necessárias novas iterações do estudo ou consultas diretas aos usuários para avaliações qualitativas. Isso porque a questão central pode não ser apenas a ordenação de um elemento específico, mas sim a interface como um todo, exigindo mais investigações para direcionar o desenvolvimento. Entretanto, dado o escopo desta pesquisa, novas iterações ou coletas de opinião dos usuários não foram realizadas.

Assim, embora não tenha sido identificada uma diferença estatisticamente significativa na qualidade em uso entre as versões, o processo contribuiu para uma tomada de decisão baseada em dados. A experimentação permitiu que os envolvidos compreendessem o comportamento do produto em uso real, possibilitando decisões fundamentadas em evidências concretas, e não apenas em suposições.
